\subsection{New-source confirmation requires memory as well as rate controls}\label{sec:newresults}
The new-source results qualify a rate-only explanation (Table~\ref{tab:newconfirm}).
At $h=24$, the CQR state gain on air quality falls from 10.59\% to 0.23\% under the
wide grid and to zero once zero/small rates are admitted. On weather it falls from
9.21\% to 6.11\% and 6.08\%, respectively: smaller rates alone leave a substantial
difference. Joint window selection then reduces it to 0.05\%. Thus the effect of
state structure cannot be assessed without its interaction with empirical score memory.

\begin{table}[!htbp]\centering\small
\caption{Locked new-source confirmation, $h=24$. Mean paired issued/shared interval-score
gain (\%), four pairs per source. Each column uses its predefined validation selector.
The final column adds window selection; it does not extend the rate grid again.}
\label{tab:newconfirm}\begin{tabular}{llrrrr}
\toprule
Source & Score & Original & Wide & Zero/small & Window \\
\midrule
Air quality & Residual & 15.04 & 3.08 & 0.00 & 2.16 \\
Air quality & CQR & 10.59 & 0.23 & 0.00 & 0.00 \\
Weather & Residual & 7.82 & 4.32 & 4.21 & 3.76 \\
Weather & CQR & 9.21 & 6.11 & 6.08 & 0.05 \\
\bottomrule
\end{tabular}

\end{table}

The same controls do not affect every score identically. Residual-score state gains
remain 2.16\% on air quality and 3.76\% on weather under joint selection. On air quality,
zero/small rates at fixed window remove the residual state difference, but allowing
window selection brings a difference back. Nested validation searches need not produce
monotonic test improvement, and the result is not an intrinsic advantage of one state.
The per-pair log contrasts preserve these interactions without dividing group averages.

For CQR, six of eight issued and seven of eight shared configurations select zero
learning rate. Both states select $W=900$ in seven of eight cases, so the memory search
also reaches a boundary: we do not claim global optimality. Directly against separately
window-tuned rolling CQR, the mean score gain is zero on air quality and 0.04\% for
issued ACI on weather ($-0.01\%$ for shared). Weather mean coverage is 89.43\% for
issued ACI versus 89.51\% for rolling. Small score differences therefore do not establish
equivalence or resolve coverage. Complete rate/window curves remain in the artifact.

Selector stability further limits interpretation. At $h=24$, early and late validation
select the identical rate/window in only three of eight issued CQR pairs and four of
eight shared CQR pairs; no residual pair has exact agreement for either state. This is
a sensitivity result, not evidence of test leakage. The primary selector remains the
full pre-test validation block specified before acquisition.

\subsection{Full configured AcMCP on the new weather subset}\label{sec:acmcpfull}
All four weather archives complete the author's 24-horizon AcMCP-CSS recursion.
The cached implementation reproduces all 96 trajectories; 60 sampled fitted scorecasts
are independently refitted. Table~\ref{tab:acmcpfull} includes every predefined pair.
At $h=24$, AcMCP-CSS has a mean paired score ratio of 1.147 relative to selected-window
rolling CQR, with 87.54\% coverage. The no-scorecast control scores better here but
has slightly lower coverage. This is a result for one declared fitting option and
four dependent weather pairs, not a general ranking of AcMCP or a controller-only
attribution. It shows that adding a faithful autocorrelation-aware comparator does
not reverse the near-rolling CQR finding on this subset. The entire grid and scorecast
success counts remain available, including initialization-related missing scorecasts.

\begin{table}[!htbp]\centering\small
\caption{Complete-configuration weather comparison, four pairs per horizon. IS ratio
averages each method's score divided by its matched validation-selected rolling CQR
score. Coverage is percent, Local is absolute coverage error in percentage points.
Ratios do not remove differences in native versus signed-error information.}
\label{tab:acmcpfull}\begin{tabular}{rlrrr}
\toprule
Horizon & Method & IS ratio & Coverage & Local \\
\midrule
1 & Rolling CQR & 1.000 & 89.37 & 3.15 \\
1 & Issued CQR & 0.996 & 89.44 & 1.43 \\
1 & AcMCP-CSS & 1.024 & 89.23 & 2.03 \\
1 & PID, no scorecast & 1.026 & 89.20 & 1.99 \\
6 & Rolling CQR & 1.000 & 90.27 & 2.91 \\
6 & Issued CQR & 1.002 & 90.44 & 2.91 \\
6 & AcMCP-CSS & 1.087 & 88.65 & 3.03 \\
6 & PID, no scorecast & 1.084 & 88.39 & 3.23 \\
24 & Rolling CQR & 1.000 & 89.51 & 6.03 \\
24 & Issued CQR & 1.000 & 89.43 & 6.01 \\
24 & AcMCP-CSS & 1.147 & 87.54 & 7.49 \\
24 & PID, no scorecast & 1.056 & 87.18 & 7.31 \\
\bottomrule
\end{tabular}

\end{table}
